\paragraph*{Confidence statements.} Beside each score, PRISM states how sure it is in three ways. It gives a set of
possible outcomes, says whether the imaging and clinical-history models agree, and reports how well the models
performed on comparable scans.

The set lists the outcomes PRISM cannot rule out at the scan (progression, no progression, or both). It is chosen so
that it includes the true outcome for 80\% of scans. It comes from split-conformal prediction
\citep{angelopoulos2023}. This method holds out calibration patients that the model was not fitted on, and uses their
scores to measure how far a score can lie from the true outcome. Write $\hat p(1) = \hat p$ for a scan's score and
$\hat p(0) = 1 - \hat p$. For each calibration scan $i$ with outcome $y_i$, the nonconformity score
$s_i = 1 - \hat p_i(y_i)$ is the probability that the model did not give to the true outcome. A patient contributes
several scans, so each calibration scan is weighted by one over the number of that patient's calibration scans, and
every patient counts once. With $\alpha = 0.2$ and $n$ calibration patients, the threshold $\hat q$ is the smallest
$s_i$ whose cumulative weight, taking the $s_i$ in increasing order, reaches $(1 - \alpha)(n + 1)$. If no $s_i$ reaches
it, $\hat q$ is infinite and every set holds both outcomes. A new scan receives every outcome to which the model gave
enough probability,
\begin{equation}
C(\hat p) = \{\, y \in \{0, 1\} : 1 - \hat p(y) \le \hat q \,\}.
\end{equation}
A score near 0 or 1 gives a single outcome, and a score near 0.5 gives both. When $\hat q < 0.5$, a score between
$\hat q$ and $1 - \hat q$ gives an empty set.

The calibration patients must be kept apart from the patients the model is fitted on. As planned, each of the three
models was therefore refitted on 70\% of its training patients and calibrated on the other 30\%. At MU this was done
inside each training fold of the cross-validation. For the external cohorts it was done once, on all MU training
patients, and the resulting models and thresholds were applied without retraining. The set of the combined score is
the primary one.

The 80\% is a rate over scans. Suppose a new patient is exchangeable with the calibration patients, which means that
swapping them would not change the joint distribution of their data. Then about 80\% of their sets contain the true
outcome on average over scans, but not necessarily within a subgroup \citep{romano2020,barber2021limits}. Weighting by
patient makes even this average approximate. The rate carries over to an external cohort only if its scans are
exchangeable with MU's \citep{mehrtens2023}, so we also calibrated within groups of comparable scans.

% Sources: METHODS_PLAN.md M2; Supplementary Section S10.
Group-conditional (Mondrian) sets compute $\hat q$ separately within each group of scans, so that the 80\% target
applies to each group and not only to the cohort as a whole \citep{vovk2022book}. A group combines the time window
after radiotherapy, defined above, with the state of the enhancing tumour. The tumour can be absent (below 0.1\,cc),
present with no previous scan to compare, increased since the previous scan by at least 1\,cc and 40\%, or not
increased. Most groups hold few calibration patients, and small calibration sets give overly wide sets
\citep{johansson2015}. Each scan is therefore calibrated at the finest of three levels that holds at least 10
calibration patients. These levels are its group, its enhancement state over all time windows, and all calibration
scans (the planned set). Clustered conformal prediction gives classes with too few calibration examples the
all-scan threshold in the same way \citep{ding2023}. LUMIERE records no radiotherapy dates, so its scans start at the
second level. Calibrating within groups can raise coverage in a group just by making sets larger. To separate the two
effects, we compared it at MU with planned-type sets at the $\alpha$ (0.175) that gives the same share of sets holding
both outcomes (exploratory).

The agreement statement says whether the clinical-history and imaging models point the same way at the scan. Each
model's score is compared with the median of that model's scores over the cohort's evaluated scans. The median uses no
outcomes, and at MU it is taken in each cross-validation repeat. The two models agree when both scores lie above
their medians or both lie below. One model sees no images and the other sees no dates, so disagreement means the
treatment course and the tumour point in different directions.

The reliability card, after model cards \citep{mitchell2019}, looks up how well the models have performed on scans in
the same group as this one. For each cohort separately, it gives how many scans with and without progression in that
group the models were tested on. At MU these are out-of-fold scans, and elsewhere they are scans scored without
retraining. With at least 10 of each, the group is in the tested tier, and the card also gives the three models' AUCs
with 95\% CIs from 1,000 bootstrap samples of patients. With fewer, the tier is few cases and the card gives only the
counts. With no scans, the tier is not tested. The record carries the tier of the scan's own cohort. The card fits
nothing. It reads results that have already been measured, including the scan it describes.
\todo{SUPPLEMENT: Recomputed without the patient's own scans, the card changed the tier of one MU scan (supplement,
``Reliability card without the described patient'', sec:S25).} The agreement and tier rules were set after the first
results at the three cohorts had been seen (exploratory).
